\documentclass[14pt,a4paper,oneside]{extbook}
\usepackage[utf8]{vietnam}
\usepackage{amsmath,amssymb,mathrsfs,fancyhdr,enumerate,multirow,makecell}
\usepackage{tkz-euclide,tikz-3dplot,tikz,tkz-tab}
\usetikzlibrary{arrows.meta,calc,intersections,angles,quotes,backgrounds,shapes.geometric,patterns,through,shadows,positioning}
\usepackage{yhmath}
\usepackage[top=0.7cm, bottom=0.7cm, left=1.0cm, right=1.0cm, headheight=16pt, headsep=3mm, footskip=5mm]{geometry}
\usepackage[hidelinks,unicode]{hyperref}
\usepackage{qrcode}
\usepackage[many]{tcolorbox}
\usepackage{fontawesome}

\pagestyle{fancy}
\fancyhf{}
\fancyhead[L]{\footnotesize\bfseries\color{teal!85!black}\faBook~SỔ TAY KIẾN THỨC TOÁN 5}
\fancyhead[R]{\footnotesize\bfseries\color{blue!85!black}Lớp Toán Thầy Lĩnh -- loptoanthaylinh.vn}
\fancyfoot[C]{\footnotesize\bfseries\color{gray!80!black}-- Sổ tay kiến thức: Bài 4. Phân số thập phân --}
\renewcommand{\headrulewidth}{0.4pt}

\begin{document}

% --- BANNER TIÊU ĐỀ BÀI HỌC ---
\noindent
\begin{tcolorbox}[
  enhanced,
  colback=teal!8!white,
  colframe=teal!80!black,
  arc=3mm,
  boxrule=1.1pt,
  width=\linewidth,
  halign=center,
  top=2.2mm, bottom=2.2mm,
  left=3mm, right=3mm,
  drop fuzzy shadow
]
  {\large\bfseries\color{teal!95!black}\faBookmark\ BÀI 04. PHÂN SỐ THẬP PHÂN}\\[2pt]
  {\scriptsize\color{black!75}\itshape Định nghĩa, cách nhận biết, chuyển đổi phân số về phân số thập phân trên 1 trang A4}
\end{tcolorbox}

\vspace{3.2mm}
% --- KHỐI 1: ĐỊNH NGHĨA & DẤU HIỆU ---
\noindent
\begin{minipage}[t]{0.56\textwidth}
  \begin{tcolorbox}[
    enhanced,
    title={\footnotesize\bfseries\color{white}\faKey\ 1. Khái niệm phân số thập phân},
    colback=blue!2!white,
    colframe=blue!75!black,
    arc=2.5mm, boxrule=0.8pt, top=1.8mm, bottom=1.8mm, left=2.2mm, right=2.2mm
  ]
    \footnotesize
    $\bullet$ \textbf{Định nghĩa:} Phân số thập phân là phân số có \textbf{mẫu số} là $10; 100; 1\,000; 10\,000; \dots$\\
    \textit{VD:} $\dfrac{3}{10}; \dfrac{47}{100}; \dfrac{215}{1000}; \dfrac{9}{10\,000}$.\\[2pt]
    $\bullet$ \textbf{Lưu ý quan trọng:} Các phân số có mẫu là $20, 50, 200\dots$ \textbf{chưa phải} là phân số thập phân, nhưng có thể chuyển thành phân số thập phân.
  \end{tcolorbox}
\end{minipage}%
\hfill
\begin{minipage}[t]{0.42\textwidth}
  \begin{tcolorbox}[
    enhanced,
    title={\footnotesize\bfseries\color{white}\faExchange\ 2. Cặp số tạo ra $10, 100, 1000$},
    colback=teal!2!white,
    colframe=teal!75!black,
    arc=2.5mm, boxrule=0.8pt, top=1.8mm, bottom=1.8mm, left=2mm, right=2mm,
    halign=center
  ]
    \footnotesize
    $2 \times 5 = 10$\\[2pt]
    $4 \times 25 = 100$ \quad $20 \times 5 = 100$\\[2pt]
    $8 \times 125 = 1\,000$ \quad $50 \times 20 = 1\,000$\\[3pt]
    {\tiny\color{teal!90!black}\bfseries Dùng các cặp số này để nhân cả tử và mẫu}
  \end{tcolorbox}
\end{minipage}

\vspace{3mm}

% --- KHỐI 2: PHƯƠNG PHÁP CHUYỂN ĐỔI ---
\noindent
\begin{tcolorbox}[
  enhanced,
  title={\footnotesize\bfseries\color{white}\faWrench\ 3. Cách chuyển một phân số thành phân số thập phân},
  colback=purple!2!white,
  colframe=purple!70!black,
  arc=2.5mm, boxrule=0.8pt, top=1.8mm, bottom=1.8mm, left=2.5mm, right=2.5mm
]
  \footnotesize
  \textbf{Cách 1: Nhân cả tử và mẫu với một số thích hợp} (khi mẫu số là ước của $10, 100, 1\,000\dots$):\\
  \quad $\triangleright$ $\dfrac{3}{5} = \dfrac{3 \times 2}{5 \times 2} = \dfrac{6}{10}$ \qquad $\triangleright$ $\dfrac{7}{25} = \dfrac{7 \times 4}{25 \times 4} = \dfrac{28}{100}$ \qquad $\triangleright$ $\dfrac{9}{125} = \dfrac{9 \times 8}{125 \times 8} = \dfrac{72}{1\,000}$.\\[2pt]
  \textbf{Cách 2: Chia cả tử và mẫu cho một số thích hợp} (khi cả tử và mẫu cùng chia hết để ra mẫu $10, 100\dots$):\\
  \quad $\triangleright$ $\dfrac{21}{30} = \dfrac{21 : 3}{30 : 3} = \dfrac{7}{10}$ \qquad $\triangleright$ $\dfrac{64}{800} = \dfrac{64 : 8}{800 : 8} = \dfrac{8}{100}$.\\[2pt]
  \textbf{Cách 3: Rút gọn trước rồi mới nhân} (nếu mẫu số ban đầu không trực tiếp nhân/chia ra tròn lũy thừa của $10$):\\
  \quad $\triangleright$ $\dfrac{12}{15} = \dfrac{12 : 3}{15 : 3} = \dfrac{4}{5} = \dfrac{4 \times 2}{5 \times 2} = \dfrac{8}{10}$.
\end{tcolorbox}

\vspace{3mm}

% --- KHỐI 3: SƠ ĐỒ TƯ DUY ---
\noindent
\begin{tcolorbox}[
  enhanced,
  title={\footnotesize\bfseries\color{white}\faSitemap\ Sơ đồ tư duy: Nhận biết \& Chuyển đổi phân số thập phân},
  colback=yellow!2!white,
  colframe=teal!85!black,
  arc=3mm, boxrule=0.9pt, top=1.8mm, bottom=1.8mm, left=1.5mm, right=1.5mm,
  halign=center
]
  \begin{tikzpicture}[
    scale=0.92,
    every node/.style={transform shape},
    root/.style={rectangle, rounded corners=5pt, fill=teal!85!black, text=white, font=\bfseries\footnotesize, inner sep=5pt, align=center},
    b1/.style={rectangle, rounded corners=4pt, fill=blue!12, draw=blue!75!black, text=blue!90!black, font=\bfseries\scriptsize, inner sep=4.5pt, align=center},
    b2/.style={rectangle, rounded corners=4pt, fill=purple!12, draw=purple!75!black, text=purple!90!black, font=\bfseries\scriptsize, inner sep=4.5pt, align=center},
    b3/.style={rectangle, rounded corners=4pt, fill=orange!12, draw=orange!75!black, text=orange!90!black, font=\bfseries\scriptsize, inner sep=4.5pt, align=center}
  ]
    \node[root] (core) at (0, 0.85) {\faBook\ PHÂN SỐ THẬP PHÂN};
    \node[b1] (n1) at (-5.2, -0.45) {\textbf{ĐẶC ĐIỂM NHẬN BIẾT}\\[2pt]Mẫu số là $10; 100; 1\,000\dots$};
    \node[b2] (n2) at (0, -0.65) {\textbf{Ý NGHĨA}\\[2pt]Bước đệm sang Số thập phân};
    \node[b3] (n3) at (5.2, -0.45) {\textbf{KỸ THUẬT CHUYỂN ĐỔI}\\[2pt]Quy đồng về mẫu $10^n$};

    \draw[teal!85!black, thick, -{Stealth[length=2mm]}] (core) -| (n1);
    \draw[teal!85!black, thick, -{Stealth[length=2mm]}] (core) -- (n2);
    \draw[teal!85!black, thick, -{Stealth[length=2mm]}] (core) -| (n3);
  \end{tikzpicture}
\end{tcolorbox}

\vspace{3mm}

% --- KHỐI 4: ĐIỀU KIỆN ĐẶC BIỆT & BẪY TRẮC NGHIỆM ---
\noindent
\begin{minipage}[t]{0.54\textwidth}
  \begin{tcolorbox}[
    enhanced,
    title={\footnotesize\bfseries\color{white}\faLightbulbO\ Khi nào một phân số chuyển được?},
    colback=green!3!white,
    colframe=green!60!black,
    arc=2.5mm, boxrule=0.8pt, top=1.8mm, bottom=1.8mm, left=2.2mm, right=2.2mm
  ]
    \footnotesize
    Một phân số tối giản \textbf{chuyển được} thành phân số thập phân khi mẫu số của nó chỉ chứa thừa số nguyên tố \textbf{2} hoặc \textbf{5} (hoặc cả hai).\\
    \textit{VD:} $\dfrac{1}{8} = \dfrac{1}{2 \times 2 \times 2}$ chuyển được; còn $\dfrac{1}{3}, \dfrac{2}{7}$ \textbf{không} chuyển được.
  \end{tcolorbox}
\end{minipage}%
\hfill
\begin{minipage}[t]{0.44\textwidth}
  \begin{tcolorbox}[
    enhanced,
    title={\footnotesize\bfseries\color{white}\faExclamationTriangle\ Mẹo tránh bẫy trắc nghiệm},
    colback=red!2!white,
    colframe=red!70!black,
    arc=2.5mm, boxrule=0.8pt, top=1.8mm, bottom=1.8mm, left=2.2mm, right=2.2mm
  ]
    \footnotesize
    $\bullet$ Đề bài hỏi: ``Trong các phân số sau, phân số nào là phân số thập phân?'': Chỉ chọn phân số có mẫu \textbf{chính xác} là $10, 100, 1000\dots$\\
    $\bullet$ Không được chọn các phân số như $\dfrac{15}{20}$ (phải chuyển đổi mới thành PSTP).
  \end{tcolorbox}
\end{minipage}

\end{document}
